% !TEX root = ../main.tex
\section{Hyperintensionality and Its Semantic Consequences}
\label{sec:hyperintensionality}

\subsection{From RE-failure to hyperintensionality}

An operator is \emph{hyperintensional} iff it can distinguish between logically
equivalent contents. Section~\ref{sec:re-failure} establishes that $\Out_M$ is
hyperintensional in exactly this sense. Via Fact~\ref{fact:kripke}, the
consequence is worth stating baldly:

\begin{fact}
\label{fact:no-kripke}
$\Out_M$ has no Kripke semantics. There is no accessibility relation
$R$---however gerrymandered, however computationally interpreted---such that
$\Out_M \phi$ holds iff $\phi$ is true at all $R$-accessible worlds.
\end{fact}

This forecloses a tempting research programme: modelling machine ``belief'' by
weakening frame conditions on a Kripke structure, dropping factivity to get
belief rather than knowledge, dropping introspection, and so on. No such
weakening helps, because RE, N, and K are frame-independent. The failure is not
at the level of \emph{which} normal logic $\Out_M$ satisfies. It is at the level
of normality itself, and indeed of classicality.

By Fact~\ref{fact:neighborhood}, it also forecloses the obvious next move. Bare
neighborhood semantics is the standard home for non-normal operators, and it is
frequently reached for on that basis, but it validates RE in every frame. An
operator that fails RE cannot be modelled in it without modification. This is
worth emphasising because the two claims---non-normality and
hyperintensionality---are easily run together, and only the first is what
neighborhood semantics was built to accommodate.

Two frameworks remain.

\subsection{Hyperintensionalised neighborhood semantics}
\label{sec:neighborhood}

In a minimal model, $\Out_M \phi$ holds at a point $w$ iff
$\lVert \phi \rVert \in N(w)$. Reading points as prompt-contexts rather than
worlds gives the neighborhood function a natural interpretation: $N(\pi)$ is the
set of contents the model is disposed to assert in context $\pi$, and the
framework imposes no closure conditions whatever on that set---no supersets, no
intersections, no unit. That models the failures of K, M, and D neatly. The
asserted set is whatever the forward pass delivers: a bare list.

To model RE-failure the framework must be hyperintensionalised. Let
neighborhoods be sets of \emph{formulations}---or of other finer-grained content
surrogates: structured propositions, guises, strings-under-interpretation---rather
than sets of worlds. This is a disciplined and well-precedented move; the
awareness structures of \citet{fagin1987} are a special case, and were introduced
for precisely the resource-bounded, presentation-sensitive phenomena at issue
here. It is, I suggest, the formally honest picture: the model's assertion state
at a context is a set of formulations, closed under nothing.

\subsection{Impossible-worlds semantics}

Alternatively \citep{rantala1982,nolan1997}, enrich the model with impossible
worlds at which formulas are evaluated by direct assignment, free of logical
closure. $\Out_M \phi$ holds at $w$ iff $\phi$ is true at all worlds in the
relevant accessibility class, where that class includes impossible worlds. Since
equivalent formulas can diverge at impossible worlds, RE fails; since an
impossible world can verify a conditional and its antecedent without its
consequent, K fails; and so on. Rantala's original motivation was logical
omniscience---modelling an agent that fails to see the consequences of what it
accepts.

Two caveats. The approach is not a free lunch: \citet{bjerring2011} proves an
impossibility result for impossible-worlds models of non-omniscient agents, and
anyone adopting the framework for this purpose owes a response to it. And the
frameworks are, for present purposes, intertranslatable; the choice is one of
illumination rather than expressive power.

I favour the hyperintensional neighborhood picture, for a reason that is
philosophical rather than technical. Impossible-worlds semantics retains the
\emph{shape} of an intentional attitude: the agent is still depicted as ruling
worlds in and out, with a coarser sieve that admits logically impossible
scenarios. That is apt for a bounded reasoner---a system trying to rule out
possibilities but lacking the resources to rule out all the impossible ones. It
subtly flatters $\Out_M$ by depicting $M$ as a possibility-manager. The
neighborhood-of-formulations picture makes no such concession. It depicts the
assertion state as what it is: a closure-free set of endorsed strings indexed to
a context, with nothing in the semantics suggesting that $M$ is representing ways
the world might be. Since my thesis is precisely that $\Out_M$ is not an
attitude toward ways the world might be, the neighborhood formulation is the
perspicuous one. Readers who prefer impossible worlds may translate throughout;
nothing downstream depends on the choice, only on the shared fact that both
frameworks are non-normal and the operator hyperintensional.

\subsection{The structural argument}
\label{sec:structural-argument}

The pieces assemble:

\begin{enumerate}
  \item Knowledge, in the received formal sense, is a normal and factive
        operator; belief, in the received formal sense, is a normal operator
        (standardly KD45). In particular both validate RE, N, and K at the level
        of competence.
  \item $\Out_M$ fails RE, K, M, D, 4, and 5 constitutively---as consequences of
        the architecture, not as performance deviations from a closure-respecting
        competence (Sections~\ref{sec:k-failure}--\ref{sec:re-failure}; the
        competence/performance premise is defended in
        Section~\ref{sec:human-parallel}).
  \item Therefore $\Out_M$ is neither a knowledge operator nor a belief operator
        in the received formal sense.
\end{enumerate}

The step from (3) to the system-level claim---that language models are neither
epistemic nor doxastic systems---requires one further premise, and it is
important not to let it pass unnoticed. $\Out_M$ is an operator over
\emph{outputs}; $K$ and $B$ are operators over \emph{attitudes}. Showing that a
system's output logic is non-classical does not by itself show that the system
has no attitudes, and the inference would prove far too much if it did: human
assertion fails RE, K, M, D, 4, and 5 just as thoroughly, and humans have
beliefs.

What licenses the step is the claim flagged in Section~\ref{sec:defining-o}:
that there is no $\pi$-invariant state governing assertion for the prompt-form
to elicit. If that holds, then $\Out_M$ is not the noisy expression of an
underlying attitude---there is no underlying attitude for it to express---and the
system-level conclusion follows. If it fails, the conclusion is confined to the
logic of machine assertion, which is a weaker but still substantive result.
Section~\ref{sec:internal} argues for the premise, and does so against the
literature most directly opposed to it.

The argument is deliberately conditional on the received formalisms, and
Section~\ref{sec:deflationism} confronts the resulting deflationary worry. But
the conditional conclusion already has teeth. Whatever machine assertion is, the
modelling apparatus built for knowledge and belief---Kripke structures,
partition models, Bayesian refinements presupposing propositional
contents---does not apply to it. Something else is needed, and
Sections~\ref{sec:grounding}--\ref{sec:subdoxastic} supply it.
